\documentclass[10pt,a4paper]{amsart}
\usepackage{amsmath,amssymb,amsthm,graphicx}
\usepackage{verbatim}
\newenvironment{reference}{\comment}{\endcomment}
%\newenvironment{reference}{}{}
\newenvironment{slogan}{\comment}{\endcomment}
\newenvironment{history}{\comment}{\endcomment}

% For commutative diagrams we use Xy-pic
\usepackage[all]{xy}

% We use 2cell for 2-commutative diagrams.
\xyoption{2cell}
% Theorem environments.
%
\theoremstyle{plain}
\newtheorem{theorem}[subsection]{Theorem}
\newtheorem{proposition}[subsection]{Proposition}
\newtheorem{lemma}[subsection]{Lemma}

\theoremstyle{definition}
\newtheorem{definition}[subsection]{Definition}
\newtheorem{example}[subsection]{Example}
\newtheorem{exercise}[subsection]{Exercise}
\newtheorem{situation}[subsection]{Situation}

\theoremstyle{remark}
\newtheorem{remark}[subsection]{Remark}
\newtheorem{remarks}[subsection]{Remarks}

\numberwithin{equation}{subsection}

% Macros
%
\def\lim{\mathop{\mathrm{lim}}\nolimits}
\def\colim{\mathop{\mathrm{colim}}\nolimits}
\def\Spec{\mathop{\mathrm{Spec}}}
\def\Hom{\mathop{\mathrm{Hom}}\nolimits}
\def\Ext{\mathop{\mathrm{Ext}}\nolimits}
\def\SheafHom{\mathop{\mathcal{H}\!\mathit{om}}\nolimits}
\def\SheafExt{\mathop{\mathcal{E}\!\mathit{xt}}\nolimits}
\def\Sch{\mathit{Sch}}
\def\Mor{\mathop{\mathrm{Mor}}\nolimits}
\def\Ob{\mathop{\mathrm{Ob}}\nolimits}
\def\Sh{\mathop{\mathit{Sh}}\nolimits}
\def\NL{\mathop{N\!L}\nolimits}
\def\CH{\mathop{\mathrm{CH}}\nolimits}
\def\proetale{{pro\text{-}\acute{e}tale}}
\def\etale{{\acute{e}tale}}
\def\QCoh{\mathit{QCoh}}
\def\Ker{\mathop{\mathrm{Ker}}}
\def\Im{\mathop{\mathrm{Im}}}
\def\Coker{\mathop{\mathrm{Coker}}}
\def\Coim{\mathop{\mathrm{Coim}}}

% Boxtimes
%
\DeclareMathSymbol{\boxtimes}{\mathbin}{AMSa}{"02}

%
% Macros for moduli stacks/spaces
%
\def\QCohstack{\mathcal{QC}\!\mathit{oh}}
\def\Cohstack{\mathcal{C}\!\mathit{oh}}
\def\Spacesstack{\mathcal{S}\!\mathit{paces}}
\def\Quotfunctor{\mathrm{Quot}}
\def\Hilbfunctor{\mathrm{Hilb}}
\def\Curvesstack{\mathcal{C}\!\mathit{urves}}
\def\Polarizedstack{\mathcal{P}\!\mathit{olarized}}
\def\Complexesstack{\mathcal{C}\!\mathit{omplexes}}
% \Pic is the operator that assigns to X its picard group, usage \Pic(X)
% \Picardstack_{X/B} denotes the Picard stack of X over B
% \Picardfunctor_{X/B} denotes the Picard functor of X over B
\def\Pic{\mathop{\mathrm{Pic}}\nolimits}
\def\Picardstack{\mathcal{P}\!\mathit{ic}}
\def\Picardfunctor{\mathrm{Pic}}
\def\Deformationcategory{\mathcal{D}\!\mathit{ef}}
\usepackage{cleveref}
\crefname{originallemma}{original lemma}{original lemmas}
\Crefname{originallemma}{Original lemma}{Original lemmas}
\makeatletter
\newlabel{lemma-localize-morphism-topoi}{{04H1}{source}{}{sourcepointer-1}{}}
\newlabel{lemma-localize-morphism-topoi@cref}{{[originallemma][0][]04H1}{[1][0][]source}{}{}{}}
\makeatother
\crefname{originallemma}{original lemma}{original lemmas}
\Crefname{originallemma}{Original lemma}{Original lemmas}
\makeatletter
\newlabel{lemma-relocalize-topos}{{04IR}{source}{}{sourcepointer-2}{}}
\newlabel{lemma-relocalize-topos@cref}{{[originallemma][0][]04IR}{[1][0][]source}{}{}{}}
\makeatother
\title{The localization of a morphism of topoi along a morphism of sheaves}
\author{The Stacks Project Authors}
\date{stacks-project:sites.tex:local-snapshot-20261012; GNU FDL 1.2 or later}
\begin{document}
\maketitle
\noindent\textit{Editorial note.} All delivered mathematical text is the author's own, copied from the pinned Stacks Project chapter file sites-cohomology.tex; no statement or proof step was rewritten, reordered or removed.\par\medskip
\noindent\textit{Editorial note.} The proof is the author's complete argument: the commutativities are reduced to the two companion diagrams and the description of the pullback is obtained by combining the two available descriptions; the closing 'which is formal' concerns an identity of fibre products, a routine verification, and it is preserved verbatim inside the counted span.\par\medskip
\noindent\textit{Editorial note.} Every labelled reference occurring inside the counted spans is carried in assembly.external\_references as a non-load-bearing navigation pointer (same-chapter ones at their definition site, cross-chapter ones at the invocation inside the counted span), so no load-bearing mathematics is replaced by a stub. No bibliography citation occurs inside the counted spans.\par\medskip
\noindent\textit{Editorial note.} The checked chapter list for same-name collisions is the set of stacks chapters staged in this repository: chow, cohomology, divisors, etale-cohomology, flat, more-algebra, more-morphisms, morphisms, spaces-more-morphisms, varieties together with this card's chapter sites-cohomology.\par\medskip
\noindent\textit{Editorial note.} The preamble is the chapter's own retained minimal preamble (sources/stacks-sites-cohomology/retained-preamble.tex), a verbatim excerpt of the native preamble: the Xy-pic setup and the native theorem declarations with the mathematical macros.\par\medskip
\noindent\textit{Allowed standard prerequisites (editorial).}\begin{itemize}
\item The localization of a morphism of topoi. Same chapter, Lemma lemma-localize-morphism-topoi. Supplies the morphism between localized topoi whose description is used in the proof.
\item Relocalizing a localization of a topos. Same chapter, Lemma lemma-relocalize-topos. Supplies the commutativity of the second square and the description of the pullback of the second localization.
\end{itemize}\medskip
\ifdefined\hypertarget\hypertarget{sourcepointer-1}{}\fi
\noindent\textit{Source reference.} Original lemma 04H1: sources/stacks-sites/pinned-original-source.tex, source lines 7904--7904. Navigation pointer invoked at proof line 7904; the counted span contains the pointer only and the cited result is not reproduced inside it, so the pointer is registered as a non-load-bearing navigation stub in the established form for Stacks material. No bibliography citation is involved.\par\medskip
\ifdefined\hypertarget\hypertarget{sourcepointer-2}{}\fi
\noindent\textit{Source reference.} Original lemma 04IR: sources/stacks-sites/pinned-original-source.tex, source lines 7951--7951. Navigation pointer invoked at proof line 7951; the counted span contains the pointer only and the cited result is not reproduced inside it, so the pointer is registered as a non-load-bearing navigation stub in the established form for Stacks material. No bibliography citation is involved.\par\medskip
\begin{lemma}
\label{lemma-relocalize-morphism-topoi}
Let $f : \Sh(\mathcal{C}) \to \Sh(\mathcal{D})$
be a morphism of topoi.
Let $\mathcal{G} \in \Sh(\mathcal{D})$,
$\mathcal{F} \in \Sh(\mathcal{C})$
and $s : \mathcal{F} \to f^{-1}\mathcal{G}$ a morphism of sheaves.
There exists a commutative diagram of topoi
$$
\xymatrix{
\Sh(\mathcal{C})/\mathcal{F} \ar[r]_{j_\mathcal{F}} \ar[d]_{f_s} &
\Sh(\mathcal{C}) \ar[d]^f \\
\Sh(\mathcal{D})/\mathcal{G} \ar[r]^{j_\mathcal{G}} &
\Sh(\mathcal{D}).
}
$$
We have $f_s = f' \circ j_{\mathcal{F}/f^{-1}\mathcal{G}}$ where
$f' :
\Sh(\mathcal{C})/f^{-1}\mathcal{G}
\to
\Sh(\mathcal{D})/\mathcal{F}$
is as in
Lemma \ref{lemma-localize-morphism-topoi}
and
$j_{\mathcal{F}/f^{-1}\mathcal{G}} :
\Sh(\mathcal{C})/\mathcal{F}
\to
\Sh(\mathcal{C})/f^{-1}\mathcal{G}$
is as in
Lemma \ref{lemma-relocalize-topos}.
The functor $(f_s)^{-1}$ is described by the rule
$$
(f_s)^{-1}(\mathcal{H} \xrightarrow{\varphi} \mathcal{G})
=
(f^{-1}\mathcal{H} \times_{f^{-1}\varphi, f^{-1}\mathcal{G}, s} \mathcal{F}
\rightarrow \mathcal{F}).
$$
Finally, given any morphisms $b : \mathcal{G}' \to \mathcal{G}$,
$a : \mathcal{F}' \to \mathcal{F}$ and
$s' : \mathcal{F}' \to f^{-1}\mathcal{G}'$ such that
$$
\xymatrix{
\mathcal{F}' \ar[r]_-{s'} \ar[d]_a & f^{-1}\mathcal{G}' \ar[d]^{f^{-1}b} \\
\mathcal{F} \ar[r]^-s & f^{-1}\mathcal{G}
}
$$
commutes, then the diagram
$$
\xymatrix{
\Sh(\mathcal{C})/\mathcal{F}'
\ar[r]_{j_{\mathcal{F}'/\mathcal{F}}} \ar[d]_{f_{s'}} &
\Sh(\mathcal{C})/\mathcal{F} \ar[d]^{f_s} \\
\Sh(\mathcal{D})/\mathcal{G}' \ar[r]^{j_{\mathcal{G}'/\mathcal{G}}} &
\Sh(\mathcal{D})/\mathcal{G}.
}
$$
commutes.
\end{lemma}

\begin{proof}
The commutativity of the first square follows from the commutativity of
the diagram in
Lemma \ref{lemma-relocalize-topos}
and the commutativity of the diagram in
Lemma \ref{lemma-localize-morphism-topoi}.
The description of $f_s^{-1}$ follows on combining the descriptions
of $(f')^{-1}$ in
Lemma \ref{lemma-localize-morphism-topoi}
with the description of
$(j_{\mathcal{F}/f^{-1}\mathcal{G}})^{-1}$ in
Lemma \ref{lemma-relocalize-topos}.
The commutativity of the last square then follows from the
equality
$$
f^{-1}\mathcal{H} \times_{f^{-1}\mathcal{G}, s} \mathcal{F}
\times_\mathcal{F} \mathcal{F}'
=
f^{-1}(\mathcal{H} \times_\mathcal{G} \mathcal{G}')
\times_{f^{-1}\mathcal{G}', s'} \mathcal{F}'
$$
which is formal.
\end{proof}

\end{document}
